% Generated draft: requires mathematical and source review.
\documentclass[11pt]{article}
\usepackage[margin=25mm]{geometry}
\usepackage{fontspec}
\setmainfont{Latin Modern Roman}
\usepackage{amsmath,amssymb,mathtools,unicode-math}
\setmathfont{Latin Modern Math}
\usepackage{enumitem,needspace,xurl,hyperref}
\setlength{\parindent}{0pt}
\setlength{\parskip}{5pt}
\setlength{\emergencystretch}{3em}
\newcommand{\N}{\mathbb{N}}
\newcommand{\Z}{\mathbb{Z}}
\newcommand{\Q}{\mathbb{Q}}
\newcommand{\R}{\mathbb{R}}
\newcommand{\C}{\mathbb{C}}
\newcommand{\F}{\mathbb{F}}
\newcommand{\A}{\mathbb{A}}
\newcommand{\PP}{\mathbb{P}}
\newcommand{\E}{\mathbb{E}}
\newcommand{\eps}{\varepsilon}
\newcommand{\dd}{\,\mathrm{d}}
\begin{document}
\section*{UnsolvedMath ID: 2303009}
\textbf{Source problem code:} catalogue entry
\subsection{Definitions and mathematical statement}
Let $D\subset\R^n$ be a convex domain with $n\geq3$, and fix a pole $Q\in D$. Let $g_D(P,Q)$ denote its positive Dirichlet Green function for the Laplacian, with a specified normalization such as $-\Delta_Pg_D(\cdot,Q)=\delta_Q$, interpreted with the standard minimal-positive convention if $D$ is unbounded. The source uses both $g$ and $G$ without defining different functions; they are treated here as references to the same Green function, subject to verification.

For $\lambda>0$, write
\[
 \Sigma_{\lambda,Q}=\{P\in D\setminus\{Q\}:g_D(P,Q)=\lambda\}.
\]
At smooth regular points of this level surface its principal curvatures are defined relative to a specified normal. Convexity of a level surface means it bounds a convex superlevel region, with the pole included in that region. The source reports Gabriel's theorem that these positive level surfaces are convex, and describes its proof as long.

\textbf{Questions (G. E. H. Reuter).} Find Green-function inequalities for convex domains in dimensions $n\geq3$ that generalize the planar consequences of classical inequalities for convex univalent functions. Find a simpler proof of the reported convexity of positive Green level surfaces. Obtain definite inequalities for their curvatures. Investigate the conjectural assertion that a half-space gives the extreme case.

The source does not identify the planar inequalities to generalize, which principal/mean/Gaussian curvature quantities are intended, the direction or constants of a bound, or the scale and pole normalization for comparison with a half-space. The requested half-space extremality is therefore a proposed research direction rather than a completely specified inequality. No particular curvature bound is invented. A comparison would need to state what is fixed, for example pole distance or an appropriate scale, because Euclidean dilation changes curvature.

The supplied summary says that Gabriel's convexity result is established, while simpler-proof, explicit-curvature, and extremality questions remain open, and reports no full resolution located. These are supplied literature assessments, not independently checked present status.
\subsection{Short English statement}
For convex domains in at least three dimensions, seek Green-function inequalities and quantitative curvature bounds for positive level surfaces, simplify Gabriel's reported convexity proof, and test the proposed extremality of a half-space. The exact inequalities and comparison normalization are not specified.
\subsection{Sources}
\begin{itemize}
\item Local UnsolvedMath registry, permanent ID 2303009, source code catalogue entry, ``Research Problems in Function Theory --- Problem 3.9''; attribution: G. E. H. Reuter. Supplied status open, classification OPEN-TRIAGE, updated 2026-08-06.
\item Hayman, \emph{Research Problems in Function Theory} (2018), Problem 3.9, \url{https://arxiv.org/abs/1809.07200}. The record reports source-TeX extraction and access on 2026-07-29; the 2018 update was retained for status triage. Current status and rights require review; the record describes publicly available arXiv source TeX.
\item The statement attributes convexity of Green level surfaces to Gabriel but supplies no publication or theorem citation. The research summary reports remaining questions unresolved without further references.
\end{itemize}
\end{document}
